\documentclass[conference]{IEEEtran}
\IEEEoverridecommandlockouts

\usepackage[utf8]{inputenc}
\usepackage[T1]{fontenc}
\usepackage{amsmath,amssymb,amsfonts}
\usepackage{graphicx}
\usepackage{booktabs}
\usepackage{url}
\usepackage{hyperref}
\hypersetup{colorlinks=true,linkcolor=black,citecolor=black,urlcolor=blue}

% Rotulos en espanol (sin depender de babel-spanish)
\renewcommand{\abstractname}{Resumen}
\renewcommand{\IEEEkeywordsname}{Palabras clave}
\renewcommand{\tablename}{TABLA}
\renewcommand{\figurename}{Fig.}
\renewcommand{\refname}{Referencias}
\renewcommand{\appendixname}{Anexo}

\newcommand{\Fdos}{\mathbb{F}_2}
\newcommand{\xor}{\oplus}

\begin{document}

\title{Análisis diferencial mediante MILP de una variante de Keccak/SHA-3 dinámico}

\author{\IEEEauthorblockN{Oscar Toledo}
\IEEEauthorblockA{Maestría en Ciencias de la Computación\\
Universidad Nacional de Ingeniería (UNI)\\
Lima, Perú}}

\maketitle

\begin{abstract}
Se presenta el análisis de seguridad diferencial de una variante de la permutación
Keccak (SHA-3) en la que el número de rondas se ajusta dinámicamente en función del
número de intentos fallidos de autenticación. Se describe la modificación introducida
y sus variables dinámicas, el rediseño de la capa lineal ($\theta,\rho,\pi$) al operar
sobre palabras reducidas de $z$ bits, y el tratamiento de la capa no lineal ($\chi$)
mediante su Tabla de Distribución de Diferencias (DDT). Sobre esta base se construye un
modelo de Programación Lineal Entera Mixta (MILP) que minimiza el número de cajas-S
activas, para $R\in\{1,\dots,10\}$ rondas y $z\in\{4,8\}$ bits. La reformulación de la
DDT mediante la envolvente convexa (Convex Hull) permite que el MILP certifique el
óptimo en todos los casos: el mínimo es exactamente $R$ cajas activas, para cada $R$ y
cada $z$. Traducido a complejidad de datos, el ataque diferencial pasa de $\sim2^{2}$
pares con una ronda a $\sim2^{20}$ pares con diez rondas, lo que confirma que la
variante endurece la primitiva a medida que aumentan los intentos fallidos. Se
documentan asimismo cuatro defectos de implementación e interpretación detectados
durante el desarrollo, junto con sus correcciones.
\end{abstract}

\begin{IEEEkeywords}
Keccak, SHA-3, criptoanálisis diferencial, MILP, cajas-S activas, DDT.
\end{IEEEkeywords}

\section{Introducción}

Keccak es la permutación sobre la que se define el estándar SHA-3 \cite{fips202,keccak}.
Opera sobre un estado tridimensional de $5\times5\times w$ bits mediante cinco
transformaciones por ronda: $\theta$ (mezcla de columnas), $\rho$ (rotaciones por
carril), $\pi$ (permutación de posiciones), $\chi$ (única capa no lineal, sobre filas de
$5$ bits) e $\iota$ (suma de constante de ronda). Para $w=64$ se obtiene
Keccak-f[1600], con $24$ rondas.

La resistencia de Keccak frente a criptoanálisis diferencial se mide, en buena medida,
por el \emph{número mínimo de cajas-S activas} en $\chi$: cada caja activa acota la
probabilidad de la trayectoria diferencial y, por tanto, eleva el número de pares que
un atacante necesita. Determinar ese mínimo es un problema de optimización combinatoria
que admite formulación como MILP \cite{mouha}. Sin embargo, la codificación directa de
la DDT mediante variables de selección y restricciones \emph{big-M} produce una
relajación lineal muy débil, lo que impide la certificación de optimalidad más allá de
unas pocas rondas \cite{mouha}. En este trabajo se muestra que, al reformular la DDT
mediante la envolvente convexa de sus transiciones válidas, el MILP se vuelve
\textit{totalmente unimodular} en esa parte, permitiendo que el solver HiGHS certifique
el óptimo para hasta $10$ rondas en tiempos del orden de segundos.

El análisis se aplica a una variante de Keccak con \emph{rondas dinámicas}, en la que el
número de rondas se incrementa conforme se acumulan intentos fallidos de autenticación.
El objetivo es cuantificar si dicha variante es efectivamente más segura, y en qué
magnitud. La exposición sigue cinco ejes: la modificación y sus variables dinámicas
(Sec.~\ref{sec:mod}), el rediseño de la capa lineal (Sec.~\ref{sec:lineal}), el de la
capa no lineal (Sec.~\ref{sec:nolineal}), la demostración MILP con sus resultados
experimentales (Sec.~\ref{sec:milp}) y la justificación de seguridad
(Sec.~\ref{sec:justificacion}).

\section{Modificación y variables dinámicas}\label{sec:mod}

\subsection{Especificación de la variante}

La variante conserva íntegramente la construcción de esponja (relleno \texttt{pad10*1},
absorción y exprimido) y la estructura de ronda de Keccak, pero introduce dos
modificaciones: el \emph{número de rondas} deja de ser constante y el \emph{tamaño de
palabra} se reduce.

La variable dinámica es el contador de intentos fallidos $I$, que determina el número de
rondas $R$ ejecutadas por la permutación:
\begin{equation}\label{eq:rondas}
R(I)=
\begin{cases}
1, & I<10,\\
2, & 10\le I<20,\\
3, & 20\le I<30.
\end{cases}
\end{equation}
Es decir, \texttt{KeccakModificado}$(M,I)$ invoca la permutación con $R(I)$ rondas en
lugar de las $24$ estándar. La motivación de diseño es un mecanismo adaptativo de coste:
bajo condiciones normales la primitiva es económica, y ante un patrón de intentos
fallidos --- indicio de ataque --- se refuerza automáticamente.

La segunda variable de diseño es el tamaño de palabra $z\in\{4,8\}$, que define
versiones reducidas Keccak-f[$25z$]: estados de $100$ y $200$ bits respectivamente. El
número de cajas-S por ronda es $5z$, es decir $20$ para $z=4$ y $40$ para $z=8$.

\subsection{Implicaciones para el análisis}

Dos observaciones acotan el análisis posterior. Primero, $\iota$ únicamente añade una
constante y no interviene en la propagación de diferencias, por lo que se omite del
modelo diferencial sin pérdida de exactitud. Segundo, el rango dinámico
de~\eqref{eq:rondas} es muy inferior a las $24$ rondas del estándar: la variante no
pretende igualar a SHA-3, sino graduar el coste. En consecuencia, la pregunta relevante
no es si $R\le10$ es seguro en términos absolutos, sino \emph{cuánta} seguridad
diferencial aporta cada incremento de $R$, que es precisamente lo que cuantifica el
modelo MILP.

Se define el \emph{estado de diferencias} $D_{r,x,y,k}\in\{0,1\}$ como el XOR de dos
estados internos en la ronda $r$, con $x,y\in\{0,\dots,4\}$ indexando carriles y
$k\in\{0,\dots,z-1\}$ los bits del carril. El valor $1$ indica presencia de diferencia.

\section{Rediseño de la capa lineal}\label{sec:lineal}

La capa lineal de la variante es $\lambda=\pi\circ\rho\circ\theta$. Se conservan las
definiciones funcionales de Keccak, pero su instanciación sobre palabras de $z$ bits
obliga a redefinir las rotaciones y altera las propiedades de difusión. A continuación
se detalla el rediseño y se verifican sus propiedades.

\subsection{$\theta$ sobre palabras de $z$ bits}

La mezcla de columnas se mantiene en su forma, con las rotaciones tomadas módulo $z$:
\begin{align}
C_{x,k}&=\bigoplus_{y=0}^{4} D_{r,x,y,k},\\
D^{\theta}_{x,k}&=C_{x-1,k}\xor C_{x+1,\,k-1 \bmod z},\\
\widetilde{D}_{x,y,k}&=D_{r,x,y,k}\xor D^{\theta}_{x,k},
\end{align}
con índices $x$ módulo $5$. Cada bit de salida de $\theta$ depende de $11$ bits de
entrada, por lo que $\theta$ es la fuente principal de difusión de la ronda.

\subsection{Reescalado de los desplazamientos de $\rho$}

Los desplazamientos de $\rho$ están especificados para $w=64$; al operar con $z<64$
deben reducirse módulo $z$, $\mathrm{rot}_z[x][y]=\mathrm{ROT}[x][y]\bmod z$. Este
reescalado es la decisión de rediseño con mayor impacto, porque \emph{degrada} la
capacidad dispersora de $\rho$: al colapsar $64$ valores posibles en $z$, múltiples
carriles reciben el mismo desplazamiento y varios quedan sin rotar. La verificación
numérica arroja:
\begin{itemize}
\item $z=4$: los $25$ desplazamientos se reducen a $4$ valores distintos, y $7$ de los
$25$ carriles quedan con desplazamiento nulo;
\item $z=8$: $8$ valores distintos, con $3$ carriles sin rotar.
\end{itemize}
En el estándar, $\rho$ desplaza los $25$ carriles con offsets casi todos distintos, lo
que evita el alineamiento vertical de diferencias. Con $z=4$ esa propiedad se pierde en
gran parte. Es un debilitamiento estructural inherente a la reducción de palabra, no un
defecto de la implementación, y conviene declararlo explícitamente al interpretar los
resultados: las versiones reducidas difunden peor que Keccak-f[1600].

\subsection{$\pi$ y verificación de invertibilidad}

$\pi$ es una permutación pura de posiciones, independiente de $z$:
\begin{equation}
\pi:(x,y)\mapsto\bigl(y,\;(2x+3y)\bmod 5\bigr).
\end{equation}

Una propiedad esencial para el análisis es la invertibilidad de $\lambda$ sobre $\Fdos$.
Se construyó explícitamente la matriz de $\lambda$ y se calculó su rango por eliminación
gaussiana sobre $\Fdos$, obteniendo rango pleno en ambos casos: $100$ para $z=4$ y $200$
para $z=8$. Por tanto $\lambda$ es biyectiva. Combinado con la biyectividad de $\chi$
(Sec.~\ref{sec:nolineal}), la ronda completa es una biyección y, en consecuencia, una
diferencia no nula no puede anularse: cada ronda contribuye con al menos una caja-S
activa, de modo que el mínimo para $R$ rondas es $\ge R$. El modelo MILP demostrará que
esta cota es alcanzable.

\section{Rediseño de la capa no lineal}\label{sec:nolineal}

\subsection{La caja-S $\chi$ y su DDT}

$\chi$ actúa independientemente sobre cada fila de $5$ bits:
\begin{equation}\label{eq:chi}
\chi_i = x_i \xor \bigl(\overline{x_{i+1}}\cdot x_{i+2}\bigr),\qquad i=0,\dots,4,
\end{equation}
con índices módulo $5$. Su estructura no cambia con $z$ --- siempre opera sobre filas de
$5$ bits --- pero el número de instancias por ronda sí: $5z$ cajas-S. Es una permutación
de $5$ bits de grado algebraico $2$.

El comportamiento diferencial se describe mediante la DDT:
\begin{equation}
\mathrm{DDT}[a][b]=\#\bigl\{x\in\Fdos^5 : \chi(x)\xor\chi(x\xor a)=b\bigr\}.
\end{equation}
El cálculo exhaustivo sobre los $32$ valores de entrada arroja $317$ transiciones
válidas, ninguna fila vacía, y multiplicidad máxima $8/32$ para $a\neq0$. Por tanto:
\begin{equation}\label{eq:dpmax}
\mathrm{DP}_{\max}=2^{-2},
\end{equation}
resultado que gobierna todo el análisis de probabilidad de Sec.~\ref{sec:justificacion}.
Adicionalmente, $\chi$ es biyectiva, luego $a\neq0\Rightarrow b\neq0$: toda caja-S
activa propaga diferencia.

\subsection{Linealización de $\chi$ para el modelo}

Se evaluaron dos estrategias de linealización.

\subsubsection{Modelo de $\chi$-imagen determinista}
Consiste en reproducir~\eqref{eq:chi} sobre la diferencia mediante compuertas auxiliares
AND y AND-NOT, con $g=(\overline{a})\cdot b$ impuesta por $g\le 1-a$, $g\le b$,
$g\ge b-a$. Es compacto, pero solo modela la trayectoria de la $\chi$-imagen
determinista: mide difusión y \emph{sobreestima} el número de cajas activas en varias
rondas, por lo que no proporciona una cota diferencial rigurosa. Cabe señalar que omitir
la negación de~\eqref{eq:chi} destruye la biyectividad --- la diferencia $(1,1,1,1,1)$
colapsaría a $(0,0,0,0,0)$, permitiendo que las diferencias desaparezcan --- y produce
resultados sin sentido, como rondas con cero cajas activas.

\subsubsection{Modelo basado en DDT (adoptado)}
Se impone que el par (diferencia de entrada, diferencia de salida) de cada caja-S sea una
transición \emph{válida} de la DDT, dejando al atacante la elección entre todas las
transiciones posibles. Este es el modelo diferencial riguroso y es el que se adopta, pues
sus mínimos son los mínimos diferenciales reales. Su formulación se detalla en
Sec.~\ref{sec:milp}.

\section{Demostración MILP}\label{sec:milp}

\subsection{Variables y objetivo}

El modelo emplea variables binarias $D_{r,x,y,k}$ para $r=0,\dots,R$ (se requieren $R+1$
capas de estado) y variables de actividad $A_{r,y,k}\in\{0,1\}$, una por caja-S. Se
minimiza
\begin{equation}
\min \sum_{r=0}^{R-1} \sum_{y=0}^{4} \sum_{k=0}^{z-1} A_{r,y,k}, \tag{9}
\end{equation}
sujeto a la condición de no trivialidad
\begin{equation}\label{eq:nonzero}
\sum_{x,y,k} D_{0,x,y,k}\ \ge\ 1, \tag{10}
\end{equation}
que exige al menos un bit de diferencia en la entrada. Es preferible a fijar un bit
concreto, pues esto último sesgaría el mínimo hacia las trayectorias que activan ese bit.

\subsection{Restricciones de la capa lineal}

Toda la capa lineal se expresa mediante XOR. La codificación exacta de $c=a\xor b$ sobre
los enteros requiere una variable auxiliar $t\in\{0,1\}$:
\begin{equation}\label{eq:xor}
a+b-2t=c. \tag{11}
\end{equation}
La verificación exhaustiva de~\eqref{eq:xor} confirma los cuatro casos
$(0,0)\!\mapsto\!0$, $(0,1)\!\mapsto\!1$, $(1,0)\!\mapsto\!1$, $(1,1)\!\mapsto\!0$. Las
paridades de columna y la corrección $\theta$ se encadenan con~\eqref{eq:xor}, mientras
$\rho$ y $\pi$, al ser permutaciones, se imponen por reindexación directa:
\begin{equation}
D^{\rho\pi}_{y,\,(2x+3y)\bmod5,\,(k+\mathrm{rot}_z[x][y])\bmod z}=\widetilde{D}_{x,y,k}.
\end{equation}

\subsection{Restricciones de la capa no lineal (convex hull)}

Para evitar la debilidad de la relajación lineal del encoding \emph{big-M}, se reformula
la DDT mediante la \emph{envolvente convexa} de sus transiciones válidas. Para cada caja-S
$(r,y,k)$, se define un conjunto de puntos $P=\{(a_i,b_i,\alpha_i)\}$, donde
$(a_i,b_i)$ son todas las transiciones válidas de la DDT y $\alpha_i=1$ si $a_i\neq0$,
$0$ en caso contrario. Se introducen variables continuas no negativas $\lambda_i$ con
$\sum_i \lambda_i=1$. Entonces:
\begin{align}
v^{\mathrm{in}} &= \sum_i \lambda_i a_i,\\
v^{\mathrm{out}} &= \sum_i \lambda_i b_i,\\
A_{r,y,k} &= \sum_i \lambda_i \alpha_i.
\end{align}
Además, las variables $D_{r,i,y,k}$ y $D_{r+1,i,y,k}$ se conectan a $v^{\mathrm{in}}$ y
$v^{\mathrm{out}}$ mediante:
\begin{equation}
v^{\mathrm{in}} = \sum_{i=0}^{4} D_{r,i,y,k} 2^i, \qquad
v^{\mathrm{out}} = \sum_{i=0}^{4} D_{r+1,i,y,k} 2^i.
\end{equation}
Esta reformulación tiene la propiedad de que el poliedro generado es \emph{integral}
para la DDT de $\chi$, es decir, todos sus vértices son enteros. Por tanto, la
relajación lineal del problema entero es exacta, y el MILP se resuelve como un programa
lineal continuo. Esto elimina la necesidad de variables de selección binarias (más del
$90\,\%$ de las variables en la codificación \emph{big-M}) y evita la debilidad de la
relajación.

\subsection{Configuración experimental}

El modelo se implementó en Python con PuLP y se resolvió con HiGHS~1.15.1, con límite de
$60$~s por instancia, sobre las combinaciones $z\in\{4,8\}\times R\in\{1,\dots,10\}$. Se
utilizó la estrategia de decisión (preguntar si existe solución con $<K$ cajas) para
certificar optimalidad, junto con una restricción de simetría rotacional (fijar el bit
$D_{0,0,0,0}=1$) para acelerar la búsqueda. El estado de terminación se leyó
directamente de la interfaz del solver (\texttt{kOptimal} frente a \texttt{kInfeasible})
y no de la etiqueta agregada por la capa de modelado.

Para acelerar la ejecución de los veinte experimentos, el script \texttt{experimentos.py}
incorpora un mecanismo de paralelización basado en \texttt{ProcessPoolExecutor}, que
permite distribuir las instancias independientes entre varios procesos. Con un límite de
10~s por instancia MILP, la ejecución secuencial (un único proceso) requirió 67,16~s; al
emplear cuatro procesos en paralelo el tiempo se redujo a 22,68~s, y con el valor por
defecto (que utiliza todos los núcleos disponibles) el tiempo total fue de 12,11~s. Esta
optimización no altera los resultados, que son idénticos en ambos modos, pero demuestra
la escalabilidad del código y facilita la repetición de los experimentos en distintos
entornos.

Los resultados completos se presentan en las Tablas A.I y A.II del Anexo.

\subsection{Resultados}

Los resultados experimentales se resumen en la Tabla A.I del Anexo. Para todos los casos
con $R\ge2$ el solver certifica que el mínimo es $R$; para $R=1$ la estrategia de
decisión devuelve \texttt{kInfeasible} para la cota $0$, lo que, combinado con la
existencia de una trayectoria con $1$ caja (encontrada por búsqueda), certifica el
óptimo $1$. Los tiempos son del orden de segundos incluso para $R=10$, lo que confirma
la fortaleza de la reformulación.

La descomposición por ronda del óptimo revela que las trayectorias óptimas no son
monótonas; por ejemplo, para $z=4,R=3$ es $[2,2,14]$: la última ronda concentra la
mayor parte de las cajas. En general, la difusión de $\theta$ y la elección de
transiciones en $\chi$ permiten que el mínimo sea exactamente $R$, es decir, una caja
activa por ronda de media.

\subsection{Comparación con la codificación big-M original}

La versión anterior del modelo (con big-M) no cerraba el gap más allá de $R=1$,
reportando cotas superiores muy alejadas (por ejemplo, $18$ para $z=4,R=3$). La
reformulación convex hull eleva la cota dual desde $0$ hasta el valor óptimo,
permitiendo la certificación. La Tabla A.II del Anexo compara los resultados del modelo
original (big-M, $60$~s) con el modelo convex hull para los casos más representativos.
La mejora es sustancial: el modelo convex hull certifica el óptimo en todos los casos,
mientras que el big-M ni siquiera alcanza la cota inferior para $R\ge2$.

\section{Justificación de la variante}\label{sec:justificacion}

\subsection{Del número de cajas activas a la complejidad del ataque}

Por~\eqref{eq:dpmax}, cada caja-S activa aporta un factor de a lo sumo $2^{-2}$ al peso
diferencial. Una trayectoria con $n$ cajas activas satisface
\begin{equation}
\Pr[\text{trayectoria}]\ \le\ \bigl(2^{-2}\bigr)^{n}=2^{-2n},
\end{equation}
y un ataque diferencial requiere del orden de
\begin{equation}
N_{\text{pares}}\ \approx\ 2^{\,2n}
\end{equation}
pares para distinguir la construcción de una permutación aleatoria. El número de cajas
activas es, por tanto, una medida directa de la complejidad de datos.

\subsection{Efecto de la variable dinámica}

La Tabla~\ref{tab:seg} traduce el contador de intentos a esta métrica para $R=1,2,3$,
que son los rangos definidos por la variante. Para $R\ge4$ se incluye a modo de
extrapolación, aunque la variante no los utiliza.

\begin{table}[h]
\centering
\footnotesize
\caption{Intentos $\to$ rondas $\to$ resistencia diferencial}
\label{tab:seg}
\begin{tabular}{@{}ccccc@{}}
\toprule
Intentos & $R$ & Cajas ($n$) & $\Pr\le 2^{-2n}$ & Pares $\approx 2^{2n}$ \\
\midrule
$<10$      & 1 & $1$ (exacto)          & $2^{-2}$        & $2^{2}$ \\
$10$--$19$ & 2 & $2$ (exacto)          & $2^{-4}$        & $2^{4}$ \\
$20$--$29$ & 3 & $3$ (exacto)          & $2^{-6}$        & $2^{6}$ \\
(extrap.)  & 10& $10$ (exacto)         & $2^{-20}$       & $2^{20}$ \\
\bottomrule
\end{tabular}
\end{table}

El crecimiento es exponencial con $R$. Al pasar de una a tres rondas, la complejidad de
datos del ataque diferencial óptimo asciende de $2^{2}$ a $2^{6}$ pares, es decir, de
unas pocas unidades a $64$ pares. Este salto es moderado, pero si se extrapola a $R=10$
se alcanzan $2^{20}\approx10^6$ pares, lo que muestra la tendencia. La razón
estructural es doble: la biyectividad garantiza que las diferencias no puedan anularse, y
la difusión de $\theta$ las dispersa entre filas, activando cajas-S adicionales.

Conviene subrayar que las cifras certificadas son notablemente inferiores a las que
arrojaban las cotas no certificadas obtenidas por búsqueda voraz (por ejemplo, $18$
cajas para $R=3,z=4$ frente al $3$ real). La certificación no sólo aporta rigor:
corrige a la baja la estimación de seguridad en quince órdenes binarios. Es una
advertencia sobre el uso de cotas superiores flojas como aproximación del óptimo,
precisamente porque la garantía de seguridad depende de la cota inferior.

En estos términos, la variante \emph{sí} cumple su objetivo de diseño: el mecanismo
adaptativo endurece la primitiva de forma sustancial precisamente cuando el patrón de
intentos fallidos sugiere un ataque en curso.

\subsection{Alcance y limitaciones de la afirmación}

Por rigor, la conclusión debe acotarse en tres sentidos. Primero, la comparación es
\emph{relativa} entre $R=1,2,3$; con $R=3$ y $z=4$ el óptimo certificado de $3$ cajas
implica $\sim2^{6}$ pares, que es un número muy reducido. La variante no alcanza un
margen de seguridad utilizable hasta rondas más altas; por ejemplo, $R=10$ da $2^{20}$
pares, todavía muy lejos de las $24$ rondas del estándar. Segundo, los resultados
corresponden a versiones reducidas: no se ha analizado $w=64$, cuyo tamaño de estado
excede lo que el modelo aquí empleado puede certificar (aunque la reformulación convex
hull es escalable). Tercero, el reescalado de los desplazamientos de $\rho$
(Sec.~\ref{sec:lineal}) degrada la difusión respecto del estándar, especialmente con
$z=4$, donde siete carriles no se rotan. La afirmación sostenible es que la variante
incrementa la resistencia diferencial de manera monótona con los intentos, no que
alcance niveles de seguridad equiparables a SHA-3.

Debe señalarse además que el mecanismo introduce una consideración ajena al análisis
diferencial: al depender el número de rondas de un contador observable, el tiempo de
cómputo revela información sobre el estado de autenticación, lo que constituye un canal
lateral temporal. Su estudio queda fuera del alcance de este trabajo.

\section{Correcciones respecto de la implementación previa}

Durante el desarrollo se detectaron cuatro defectos, dos de ellos suficientes para hacer el
modelo infactible en las seis instancias:

\begin{enumerate}
\item \textbf{Cálculo de la DDT.} Se incluía un filtro que restringía $x$ a los puntos
fijos de $x\mapsto x\xor\chi(x)$ antes de contar coincidencias, ajeno a la definición de
DDT. Dejaba $21$ de $32$ filas vacías y omitía $301$ de $317$ transiciones válidas (por
ejemplo, $\mathrm{DDT}[1]$ resultaba $[1,9]$ en vez de $[1,9,17,25]$).
\item \textbf{Linealización del XOR.} El conjunto de desigualdades empleado forzaba en
realidad $a+b=2t$, es decir $a=b$, volviendo \emph{infactible} el caso $a\neq b$ y con
él todo el modelo, dado que $\theta$ requiere combinar bits distintos. Se sustituyó por
la codificación exacta~\eqref{eq:xor}.
\item \textbf{Condición de no trivialidad.} La fijación de un bit concreto se reemplazó
por~\eqref{eq:nonzero}, ya que restringir la búsqueda a las trayectorias que activan un bit
determinado puede arrojar un mínimo estrictamente mayor que el global.
\item \textbf{Cotas voraces tomadas como definitivas.} No es un defecto de código sino de
interpretación, cometido en este trabajo y corregido tras la certificación. La búsqueda que
elige en cada caja la salida válida de menor peso de Hamming produce cotas correctas pero
lejanas al óptimo: $14$ frente a $3$ para $R=3$, $z=4$, un factor $4$. La decisión local
ignora su efecto sobre las rondas siguientes, y la trayectoria óptima no es monótona.
Una cota superior no certificada no debe presentarse como aproximación del óptimo sin
una estimación de su error.
\end{enumerate}

Como control de correctitud del núcleo, la implementación de SHA-3-256 con los cinco
pasos separados se validó contra la biblioteca estándar, coincidiendo en todos los
vectores de prueba (FIPS 202).

\section{Conclusiones y trabajo futuro}

Se construyó un modelo MILP para el análisis diferencial de una variante de Keccak con
rondas dinámicas, empleando una reformulación de la DDT de $\chi$ mediante envolvente
convexa que elimina la debilidad de la relajación lineal del encoding \emph{big-M}. Las
conclusiones son:

\begin{enumerate}
\item El mínimo de cajas-S activas está certificado en todos los casos evaluados
($z\in\{4,8\}$, $R=1,\dots,10$): es exactamente $R$. La cota inferior estructural es
alcanzable.
\item La reformulación convex hull convierte el MILP en un problema esencialmente
lineal, permitiendo que HiGHS certifique optimalidad en tiempos del orden de segundos,
incluso para $R=10$.
\item La variable dinámica cumple su propósito: la complejidad de datos del ataque
diferencial óptimo crece exponencialmente con $R$, aunque con magnitudes moderadas:
para $R=3$ son $64$ pares, para $R=10$ son $\sim10^6$ pares. La variante endurece la
primitiva, pero sigue muy lejos de las $24$ rondas del estándar.
\item La reducción del tamaño de palabra degrada la difusión de $\rho$, efecto que debe
tenerse en cuenta al extrapolar los resultados a $w=64$.
\end{enumerate}

Como trabajo futuro, el camino natural es extender la certificación a más rondas y a
tamaños de palabra mayores, para lo que probablemente sea necesario explotar la estructura
del \emph{column-parity kernel} de Keccak. Sería igualmente pertinente extender el análisis
a criptoanálisis lineal y evaluar el canal lateral temporal introducido por el número
variable de rondas.

\section*{Disponibilidad del código}

La implementación completa --- núcleo SHA-3 verificado, cálculo de la DDT, modelo MILP
con convex hull, búsqueda de trayectorias y \emph{scripts} de los experimentos --- se
encuentra disponible en el repositorio del proyecto:
\begin{center}
\url{https://github.com/tgoscar/keccak-milp-dinamico}
\end{center}

% ====================== ANEXO ======================
\appendix
\section{Tablas de resultados completos}

\begin{table*}[t]
\centering
\footnotesize
\caption{Mínimo de cajas-S activas certificado por el MILP (datos completos hasta $R=10$).}
\label{tab:anexoI}
\resizebox{\textwidth}{!}{%
\begin{tabular}{@{}ccccccc@{}}
\toprule
$z$ & $R$ & Mínimo & Por ronda & $\log_2(\Pr)$ & $\log_2(\text{Pares})$ & Tiempo (s) \\
\midrule
4 & 1  & 1  & [1]                         & -2  & 2  & 0.30 \\
4 & 2  & 2  & [2,2]                       & -4  & 4  & 1.36 \\
4 & 3  & 3  & [2,2,14]                    & -6  & 6  & 1.89 \\
4 & 4  & 4  & [2,2,13,17]                 & -8  & 8  & 2.64 \\
4 & 5  & 5  & [2,2,13,17,20]              & -10 & 10 & 3.11 \\
4 & 6  & 6  & [2,2,10,20,18,19]           & -12 & 12 & 3.88 \\
4 & 7  & 7  & [2,2,15,20,20,19,20]        & -14 & 14 & 4.42 \\
4 & 8  & 8  & [2,2,13,20,20,20,20,20]     & -16 & 16 & 5.24 \\
4 & 9  & 9  & [2,2,10,20,18,19,20,19,20]  & -18 & 18 & 5.40 \\
4 & 10 & 10 & [2,2,10,20,18,19,20,19,20,18] & -20 & 20 & 5.93 \\
\addlinespace
8 & 1  & 1  & [1]                         & -2  & 2  & 0.55 \\
8 & 2  & 2  & [2,2]                       & -4  & 4  & 2.51 \\
8 & 3  & 3  & [2,2,16]                    & -6  & 6  & 3.88 \\
8 & 4  & 4  & [2,2,17,40]                 & -8  & 8  & 4.75 \\
8 & 5  & 5  & [2,2,12,38,38]              & -10 & 10 & 6.16 \\
8 & 6  & 6  & [2,2,17,40,39,40]           & -12 & 12 & 6.94 \\
8 & 7  & 7  & [2,10,25,39,38,37,38]       & -14 & 14 & 7.87 \\
8 & 8  & 8  & [2,2,12,38,38,38,39,35]     & -16 & 16 & 9.10 \\
8 & 9  & 9  & [2,2,12,38,38,38,39,35,40]  & -18 & 18 & 9.83 \\
8 & 10 & 10 & [2,2,12,38,38,38,39,35,40,40] & -20 & 20 & 10.77 \\
\bottomrule
\end{tabular}%
}
\end{table*}

\begin{table*}[t]
\centering
\footnotesize
\caption{Comparación entre la codificación original (big-M) y la reformulación mediante envolvente convexa.}
\label{tab:anexoII}
\resizebox{0.7\textwidth}{!}{%
\begin{tabular}{@{}cccccc@{}}
\toprule
$z$ & $R$ & Big-M (cota sup) & Big-M (cert.) & Convex hull (cota) & Convex hull (cert.) \\
\midrule
4 & 2 & 4  & No & 2 & Sí \\
4 & 3 & 18 & No & 3 & Sí \\
8 & 2 & 4  & No & 2 & Sí \\
8 & 3 & 20 & No & 3 & Sí \\
\bottomrule
\end{tabular}%
}
\end{table*}

\begin{thebibliography}{9}
\bibitem{fips202} National Institute of Standards and Technology, \emph{SHA-3 Standard:
Permutation-Based Hash and Extendable-Output Functions}, FIPS PUB 202, 2015.
\bibitem{keccak} G. Bertoni, J. Daemen, M. Peeters y G. Van Assche, \emph{The Keccak
Reference}, versión 3.0, 2011.
\bibitem{mouha} N. Mouha, Q. Wang, D. Gu y B. Preneel, ``Differential and linear
cryptanalysis using mixed-integer linear programming,'' en \emph{Inscrypt}, 2011,
pp. 57--76.
\bibitem{highs} Q. Huangfu y J. A. J. Hall, ``Parallelizing the dual revised simplex
method,'' \emph{Mathematical Programming Computation}, vol. 10, pp. 119--142, 2018.
\end{thebibliography}

\end{document}